\newtheoremstyle{Definiciones}
  {10pt}         % Espacio arriba
  {10pt}         % Espacio abajo
  {\normalfont}  % Fuente del cuerpo
  {0em}          % Sangría del cuerpo (puedes usar 20pt, 1cm, etc.)
  {\bfseries}    % Fuente del encabezado (Definición 3.1)
  {}            % Puntuación después del encabezado
  {.5em}         % Espacio después del encabezado
  {\thmname{#1}\thmnumber{ #2. }\thmnote{#3}} % Formato del encabezado

\newtheoremstyle{Proposiciones}
  {10pt}         % Espacio arriba
  {10pt}         % Espacio abajo
  {\normalfont}  % Fuente del cuerpo
  {0em}          % Sangría del cuerpo (puedes usar 20pt, 1cm, etc.)
  {\bfseries}    % Fuente del encabezado (Definición 3.1)
  {}            % Puntuación después del encabezado
  {.5em}         % Espacio después del encabezado
  {\thmname{#1}\thmnumber{ #2. }\thmnote{\normalfont#3}} % Formato del encabezado

\newtheoremstyle{Demostraciones}
  {-5pt}         % Espacio arriba
  {10pt}         % Espacio abajo
  {\normalfont}  % Fuente del cuerpo
  {0em}          % Sangría del cuerpo (puedes usar 20pt, 1cm, etc.)
  {}    % Fuente del encabezado (Definición 3.1)
  {}            % Puntuación después del encabezado
  {0em}         % Espacio después del encabezado
  {\thmname{#1}\thmnumber{ #2.}\thmnote{\normalfont(#3)}} % Formato del encabezado

\newtheoremstyle{Lemas}
  {5pt}         % Espacio arriba
  {10pt}         % Espacio abajo
  {\normalfont}  % Fuente del cuerpo
  {0em}          % Sangría del cuerpo (puedes usar 20pt, 1cm, etc.)
  {\bfseries}    % Fuente del encabezado (Definición 3.1)
  {}            % Puntuación después del encabezado
  {.3em}         % Espacio después del encabezado
  {\thmname{#1}\thmnumber{ #2.}\thmnote{\normalfont(#3)}} % Formato del encabezado

\newtheoremstyle{Observaciones}
  {5pt}         % Espacio arriba
  {10pt}        % Espacio abajo
  {\normalfont} % Fuente del cuerpo
  {}            % Sangría del cuerpo 
  {\bfseries}   % Fuente del encabezado (Definición 3.1)
  {.}           % Puntuación después del encabezado
  {.5em}           % Espacio después del encabezado
  {\thmname{#1}\thmnumber{ #2}\thmnote{ (#3)}} % Formato del encabezado

\newtheoremstyle{entorno con itemize}
  {5pt}         % Espacio arriba
  {10pt}         % Espacio abajo
  {\normalfont} % Fuente del cuerpo
  {}            % Sangría del cuerpo (puedes usar 20pt, 1cm, etc.)
  {\bfseries}   % Fuente del encabezado (Definición 3.1)
  {.}           % Puntuación después del encabezado
  {0pt}         % Espacio después del encabezado
  {\thmname{#1}\thmnumber{ #2}\thmnote{ (#3)}} % Formato del encabezado

% ==============================================================================
\theoremstyle{Definiciones}
\newtheorem{deff1}{Definición}[section]
\newenvironment{deff}[2][]
  {\begin{deff1}[#1] (#2).\leavevmode \vspace{-0.1\baselineskip}}
  {\end{deff1}}

\theoremstyle{Proposiciones}
\newtheorem{_prop}[deff1]{Proposición}
\newenvironment{prop}[1][]
  {\begin{_prop}[#1]}
  {\end{_prop}}

\newtheorem{_teo}[deff1]{Teorema}
\newenvironment{teorema}[1][]
  {\begin{_teo}[#1] \leavevmode \vspace{-0.1\baselineskip} }
  {\end{_teo}}
% ==============================================================================
\theoremstyle{Observaciones}

\newtheorem{ej}{Ejemplo}[section]
\newtheorem*{obs}{Observaciones}
\newtheorem*{obs y ejs}{Observaciones y Ejemplos}

% ==============================================================================
\theoremstyle{entorno con itemize}

\newtheorem*{_nott}{Notación}
\newenvironment{nott}
  {\begin{_nott} \leavevmode \vspace{-0.4\baselineskip} \begin{itemize}}
  {\end{itemize}\end{_nott}}

\newtheorem*{_obss}{Observaciones}
\newenvironment{obss}
  {\begin{_obss} \leavevmode \vspace{-0.4\baselineskip} \begin{itemize}}
  {\end{itemize}\end{_obss}}
  
\newtheorem{_ejs}[ej]{Ejemplo}
\newenvironment{ejs}
  {\begin{_ejs}\leavevmode \vspace{-0.2\baselineskip} \begin{itemize}}
  {\end{itemize}\end{_ejs}}

% ==============================================================================
\theoremstyle{Lemas}

\newtheorem{_lema}[deff1]{Lema}
\newenvironment{lema}[1][]
  {\begin{_lema} #1 \leavevmode }
  {\end{_lema}}
% ==============================================================================
\theoremstyle{Demostraciones}

\newtheorem*{_dem}{}

\newenvironment{dem}[2][]
  {\begin{_dem} #1\textit{Demostración} (#2). }
  {\end{_dem}}